\section*{अध्याय \ref{ch:b3:group-theory-applied} --- समूह सिद्धांत के अनुप्रयोग}

\begin{solution}{exo:b3:group-theory-applied:1}
$n_{A_1} = \frac14(5 - 1 + 3 + 1) = 2$, $n_{A_2} = \frac14(5 - 1 - 3 - 1) = 0$,
$n_{B_1} = \frac14(5 + 1 + 3 - 1) = 2$, $n_{B_2} = \frac14(5 + 1 - 3 + 1) = 1$:
$\Gamma = 2A_1 + 2B_1 + B_2$, विमा 5 का।
\end{solution}

\begin{solution}{exo:b3:group-theory-applied:2}
स्थान पर छोड़े गए परमाणुओं की गिनती जल जैसी ही: $\Gamma_{\mathrm{vib}} = 2A_1 + B_1$। $A_1$ ($z$) और
$B_1$ ($x$) \omterm{def:b3:group-theory-applied:activity}{अवरक्त-सक्रिय} हैं, और दोनों \omterm{def:b3:group-theory-applied:activity}{रामन-सक्रिय} भी ($x^2$, $xz$): हर स्पेक्ट्रम में तीन
बैंड।
\end{solution}

\begin{solution}{exo:b3:group-theory-applied:3}
$B_1\otimes B_2$: \omterm{def:b3:point-groups:representation}{अभिलक्षण} $(1, 1, -1, -1)$ $= A_2$। $A_2\otimes B_1$: $(1, -1, -1,
1) = B_2$। $C_{3v}$ में $E\otimes E$ के \omterm{def:b3:point-groups:representation}{अभिलक्षण} $(4, 1, 0)$ हैं; लघुकरण: $A_1 + A_2 +
E$।
\end{solution}

\begin{solution}{exo:b3:group-theory-applied:4}
अवरक्त: दोनों $t_2$ विधाएँ, 3019 और \qty{1306}{cm^{-1}}। रामन: चारों ($a_1$,
$e$, $t_2$ सभी द्विघाती फलनों की तरह रूपांतरित होते हैं)। \omterm{def:b3:point-groups:inversion}{प्रतिलोमन केंद्र} नहीं है, अतः
संपातन अनुमत हैं।
\end{solution}

\begin{solution}{exo:b3:group-theory-applied:5}
स्थान पर छोड़े गए परमाणु $(4, 1, 2, 4, 1, 2)$ गुणा $\chi_{xyz} = (3, 0, -1, 1, -2, 1)$ से
$\chi_{3N} = (12, 0, -2, 4, -2, 2)$ मिलता है, जिसका लघुकरण $A_1' + A_2' + 3E' + 2A_2'' +
E''$ है। $E' + A_2''$ (स्थानांतरण) और $A_2' + E''$ (घूर्णन) हटाइए: $A_1' + 2E'
+ A_2''$। अवरक्त: $2E' + A_2''$, तीन बैंड; रामन: $A_1' + 2E'$, तीन बैंड। $A_2''$
(तल से बाहर बंकन) केवल अवरक्त में, $A_1'$ (सममित तनन) केवल रामन में, और $E'$
विधाएँ दोनों में दिखती हैं।
\end{solution}

\begin{solution}{exo:b3:group-theory-applied:6}
अवरक्त: दोनों $T_{1u}$ विधाएँ, दो बैंड। रामन: $A_{1g}$, $E_g$, $T_{2g}$, तीन बैंड।
$T_{2u}$ मूक है। \ce{SF6} में \omterm{def:b3:point-groups:inversion}{प्रतिलोमन केंद्र} है: परस्पर अपवर्जन।
\end{solution}

\begin{solution}{exo:b3:group-theory-applied:7}
$E$, $C_3$, $C_3^2$ के अंतर्गत फलन $h_2$ क्रमशः $h_2$, $h_3$, $h_1$ पर जाता है (\omterm{def:b3:point-groups:representation}{अभिलक्षण} 2,
$-1$, $-1$); परावर्तनों का योगदान 0 है। $\hat P_Eh_2 \propto 2h_2 - h_3 - h_1$।
$2h_1 - h_2 - h_3$ के साथ इसका अतिव्यापन (परमाणु अतिव्यापनों की उपेक्षा करके) $-3$ है, और
बाद वाले का मानक $\sqrt6$ है; प्रक्षेप घटाने पर, $(2h_2 - h_1 - h_3) +
\frac12(2h_1 - h_2 - h_3) = \frac32(h_2 - h_3)$।
\end{solution}

\begin{solution}{exo:b3:group-theory-applied:8}
फ़ैक ($C_{3v}$, \omterm{def:b3:point-groups:representation}{अभिलक्षण} $3, 0, 1$): $A_1 + E$, दो अवरक्त बैंड। मेर ($C_{2v}$,
\omterm{def:b3:point-groups:representation}{अभिलक्षण} $3, 1, 3, 1$): $2A_1 + B_1$, तीन बैंड। \ce{M(CO)6} ($O_h$): $A_{1g} + E_g
+ T_{1u}$, एक अवरक्त बैंड ($T_{1u}$)।
\end{solution}

\begin{solution}{exo:b3:group-theory-applied:9}
\omterm{def:b3:point-groups:representation}{अभिलक्षण} $6, 0, 0, 2, 2, 0, 0, 0, 4, 2$: केवल $E$, अक्षों के अनुदिश $C_4$ और $C_2$
(स्थान पर छोड़े गए दो लिगैंड), तीनों $\sigma_h$ (चार) और छहों $\sigma_d$
(दो) लिगैंडों को स्थान पर रखते हैं। लघुकरण $A_{1g} + E_g + T_{1u}$ देता है। धातु के $s$
($a_{1g}$), $d_{z^2}$ और $d_{x^2-y^2}$ ($e_g$), तथा $p_x$, $p_y$, $p_z$ ($t_{1u}$) संयोजित होते हैं;
$t_{2g}$ अनाबंधी रहते हैं।
\end{solution}

\begin{solution}{exo:b3:group-theory-applied:10}
$n_i = \frac1h(1\cdot h\cdot\chi_i(E)) = d_i$। नियमित
निरूपण की विमा $h$ है, और साथ ही $\sum_in_id_i = \sum_id_i^2$।
\end{solution}

\begin{solution}{exo:b3:group-theory-applied:11}
$D_{2h}$ में (अणु $z$ के अनुदिश), स्थान पर छोड़े गए परमाणु $(4, 4, 0, 0, 0, 0, 4, 4)$, $\chi_{3N}
= (12, -4, 0, 0, 0, 0, 4, 4)$, जिसका लघुकरण $2A_g + 2B_{2g} + 2B_{3g} + 2B_{1u} + 2B_{2u}
+ 2B_{3u}$ है। स्थानांतरण $B_{1u} + B_{2u} + B_{3u}$ और दोनों घूर्णन
$B_{2g} + B_{3g}$ हटाइए: $2A_g + B_{2g} + B_{3g} + B_{1u} + B_{2u} + B_{3u}$, सात विधाएँ।
$D_{\infty h}$ में: $2\Sigma_g^+ + \Pi_g + \Sigma_u^+ + \Pi_u$। अवरक्त: $\Sigma_u^+$ और $\Pi_u$
(दो बैंड); रामन: $2\Sigma_g^+$ और $\Pi_g$ (तीन बैंड); कोई संपातन नहीं।
\end{solution}

\begin{solution}{exo:b3:group-theory-applied:12}
$1b_2 \to 4a_1$: $B_2$, $y$, अनुमत, तल के लंबवत। $3a_1 \to 4a_1$:
$A_1$, $z$, अनुमत, $C_2$ अक्ष के अनुदिश। $1b_2 \to 2b_1$: $A_2$, वर्जित।
$1b_1 \to 4a_1$: $B_1$, $x$, अनुमत, तल में, अक्ष के लंबवत।
\end{solution}

\begin{solution}{pb:b3:group-theory-applied:1}
\textbf{1.} सिस: दोनों CO सन्निकट शीर्षों पर; ट्रांस: सम्मुख; फ़ैक: तीनों
CO एक फलक पर; मेर: तीनों CO एक याम्योत्तर पर (दो ट्रांस, एक बीच में)।
\textbf{2.} $C_{2v}$।
\textbf{3.} $D_{4h}$।
\textbf{4.} $C_{3v}$।
\textbf{5.} $C_{2v}$।
\textbf{6.} केवल trans-\ce{ML4(CO)2}।
\textbf{7.} जो संक्रिया किसी CO को दूसरे पर ले जाती है, वह उसके सदिश को
विकर्ण से बाहर रखती है (योगदान 0); जो किसी CO को स्थान पर छोड़ती है, वह उसके आबंध सदिश को
स्वयं पर प्रतिचित्रित करती है (+1); कोई भी C--O सदिश को उलटती नहीं।
\textbf{8.} \omterm{def:b3:point-groups:representation}{अभिलक्षण} $(2, 0, 2, 0)$, जहाँ $\sigma_v(xz)$ दोनों CO का तल है: $A_1 +
B_1$।
\textbf{9.} $E$, $C_4$, $C_2$, $\sigma_v$, $\sigma_d$ के अंतर्गत \omterm{def:b3:point-groups:representation}{अभिलक्षण} 2, अन्यत्र 0:
$A_{1g} + A_{2u}$।
\textbf{10.} $(3, 0, 1)$: $A_1 + E$।
\textbf{11.} $(3, 1, 3, 1)$: $2A_1 + B_1$।
\textbf{12.} 2, 2, 3, 3: हर CO के लिए एक विमा।
\textbf{13.} सिस 2, ट्रांस 1, फ़ैक 2, मेर 3।
\textbf{14.} सिस 2, ट्रांस 1, फ़ैक 2, मेर 3।
\textbf{15.} ट्रांस: उसका अवरक्त बैंड ($A_{2u}$) और रामन बैंड ($A_{1g}$) भिन्न
विधाएँ हैं।
\textbf{16.} $A_{2u}$, $i$ के अंतर्गत चिह्न बदलता है: एक CO लंबा होता है जबकि दूसरा
छोटा, विपरीत कला में।
\textbf{17.} $E$ एक अपभ्रष्ट युग्म है: समान आवृत्ति की दो विधाएँ एक
बैंड देती हैं, साथ में $A_1$ बैंड।
\textbf{18.} एक ट्रांस-द्विप्रतिस्थापित, केंद्रसममित समावयवी।
\textbf{19.} तीन अवरक्त बैंड: मेर (फ़ैक दो देता)।
\textbf{20.} हाँ: मेर के तीन \omterm{def:b3:group-theory-applied:activity}{रामन-सक्रिय} CO विधाएँ उन्हीं आवृत्तियों पर हैं,
क्योंकि वह केंद्रसममित नहीं है; फ़ैक दो दिखाता।
\textbf{21.} $A_1$: समान-कला तनन पूर्णतः सममित है।
\textbf{22.} फ़ैक/मेर मिश्रण अधिकतम पाँच अवरक्त बैंड दिखाता, जिनमें फ़ैक
बैंड भिन्न स्थितियों पर होते; तीन स्पष्ट बैंड एक ही समावयवी से मेल खाते हैं।
${}^{13}$C NMR स्पेक्ट्रम निर्णय करेगा: मेर के लिए 2:1 अनुपात में दो कार्बोनिल
संकेत, फ़ैक के लिए एक संकेत।
\textbf{23.} \textbf{मेर समावयवी के तीन \omterm{def:b3:group-theory-applied:activity}{अवरक्त-सक्रिय} CO तनन हैं, फ़ैक
समावयवी के दो}: उत्पाद मेर है।
\end{solution}
